\section{Introduction}
\label{sec:intro}

3D Gaussian Splatting (3DGS)~\cite{kerbl2023} reconstructs scenes with explicit Gaussian primitives and renders new views in real time.
Anchor-based variants~\cite{lu2024} cut the primitive count by attaching neural Gaussians to shared anchors, and a line of codecs built on this representation~\cite{chen2024,chen2025,wang2024b} compresses scenes to a few megabytes with learned context models and entropy coding.
These codecs, however, emit a \emph{single-rate} stream: the encoder fixes one rate--quality trade-off, and the receiver gets exactly one file at one quality.
Every deployment where bandwidth is unknown or shared---streaming to mobile clients, archival with future re-use, pay-as-you-download delivery---would rather have the classical property of scalable coding~\cite{taubman2002jpeg2000}: one file whose quality grows with the bytes received, so that truncating the transfer truncates the quality, not the usability.

Recent work brings progressive transmission to 3DGS.
PCGS~\cite{chen2026pcgs} schedules the rate--distortion trade-off during training so that one model serves several rates; ProGS~\cite{tang2026progs} organizes anchors as parent-closed octree prefixes; SCAR-GS~\cite{revilla2026scar} refines residual vectors across quality layers; GoDe~\cite{disario2025gode} trains progressive level-of-detail Gaussians; Octree-GS~\cite{ren2024octree} selects levels of detail at rendering time.
Before building on them, we inspected what their decoders must receive before the first anchor attribute can be decoded.
None of the five starts from plain content: each stream begins with learned context models, hash grids, codebooks, or dictionaries.
These headers are small relative to a full scene but they are fixed costs that precede any decodable byte, and they tie every prefix to a trained model that must travel with it.

\begin{figure}[t]
\centering
\includegraphics[width=\linewidth]{figs/fig1_bitstream.pdf}
\caption{\textbf{One file, four decodable tiers.}
The bitstream stores, layer by layer, a base layer (anchor geometry plus coarsest quantization symbols) followed by refinement layers.
Cutting the file at any block boundary yields a valid stream: the decoder renders the complete scene from the prefix alone, and the full-precision tier is bit-exact with the single-rate decode.
No learned model, codebook, hash grid, or quantization-parameter table precedes the first decodable byte---all of it is recomputed from received bytes.
Annotated sizes and PSNR are measured on the UAV scene of Sec.~\ref{sec:experiments} at the tightest anchor budget ($r{=}0.50$).}
\label{fig:bitstream}
\end{figure}

Keeping the header empty is the central technical obstacle, and it has two parts.
First, entropy coding needs priors: per-element quantization steps, conditioning structure, and probability models.
The HAC++ line~\cite{chen2024,chen2025} produces these from a hash grid over anchor coordinates---a learned model in the stream by construction.
Dropping it is only sound if \emph{every} quantity the decoder needs is a function of data it has already decoded, which constrains the design of the quantization, the conditioning, and the file layout simultaneously.
Second, coarse-first decoding fights the training objective: a model optimized for full precision places quantization symbols anywhere between coarse lattice points, so the first-decoded tier absorbs a large quantization error; spending encoder bits to fix the base layer defeats the purpose of layering, so the \emph{training} has to know about the ladder.

\DCCA addresses both obstacles.
Its bitstream stores anchor geometry and a base layer of coarsest-step symbols first, followed by refinement layers of binary decisions (zero-flag, sign, magnitude bits) whose probabilities come from small static frequency tables transmitted inline.
Quantization steps are recomputed at the decoder from already-decoded structure; conditioning groups (contribution rank) and probability buckets (already-decoded coarse magnitude) are recomputed likewise.
On top of this structure we add two training-side mechanisms: a ladder-alignment penalty that pulls quantization symbols toward base-layer lattice points, and a training-time anchor budget that scales the entry size of the whole stream.
Our claims and evidence:
\begin{itemize}
\item \textbf{Zero-side-information progressive bitstream.}
A single file in which every tier is a contiguous byte prefix, decodable with no learned model, codebook, hash grid, or quantization-parameter table; the full-precision tier is verified bit-exact against the single-rate decode, and decoder-recomputed quantization steps match the encoder's on 100\% (feature, offset) and 99.999\% (scaling) of symbols (Sec.~\ref{sec:method}, \ref{sec:verify}).
\item \textbf{Ladder-aligned quantization training.}
A late-training penalty toward the base-layer lattice raises the entry tier by $+0.41$\,dB at noise-level full-precision cost in the recommended configuration, and by $+1.09$\,dB when the budget is tight; we delimit where it does not help (Sec.~\ref{sec:align-exp}).
\item \textbf{Anchor-budget axis.}
A training-time ADMM budget on anchor count moves the entry tier continuously across $6.33$--$15.88$\,MB (quality 24.25--26.15\,dB); the tightest whole progressive file (11.41\,MB) stays within 6\% of the lowest tier reported by ProGS~\cite{tang2026progs}, and its entry tier is 41\% below it (Sec.~\ref{sec:budget-exp}).
\end{itemize}
All evidence follows one protocol: paired arms that share code and seed, gate thresholds written down before launch, three seed pairs to bound noise, and numbers only from decoded bitstreams (Sec.~\ref{sec:setup}).
We also report what did not work: three attempts at learned coding inside this stream, each negative, which is the quantitative case for the static, decoder-recomputable design (Sec.~\ref{sec:negative}).
